% TOP_PROBLEM: {"id": 2562, "title": "Kourovka Notebook Problem 21.53", "category_id": 17, "category": {"id": 17, "name": "group_theory", "display_name": "Group Theory", "description": "Problems about groups, group actions, representations, and related algebraic structures.", "slug": "group-theory", "order_index": 17, "created_at": "2026-07-31T15:26:25.671Z"}, "status": "open"}
% FIRST_POSTED: null
% ENTRY_KIND: statement_only
% Imported from: batch09.tex; UnsolvedMath ID 2562
\documentclass[11pt]{article}
\usepackage[margin=25mm]{geometry}
\usepackage{fontspec}\setmainfont{Latin Modern Roman}
\usepackage{amsmath,amssymb,mathtools}
\usepackage{unicode-math}\setmathfont{Latin Modern Math}
\usepackage{xurl,hyperref}
\hypersetup{colorlinks=true,urlcolor=blue}
\setcounter{secnumdepth}{0}\setlength{\parindent}{0pt}\setlength{\parskip}{5pt}\setlength{\emergencystretch}{3em}
\newcommand{\sref}[2]{\hyperlink{source:#1:#2}{[S#2]}}
\newcommand{\cataloguescope}[1]{\paragraph{Recorded target.}#1}
\begin{document}
% UnsolvedMath ID 2562; source catalogue code KOU-21.53
\setcounter{section}{2561}
\section{Kourovka Notebook Problem 21.53}\label{2562}
{\small\sffamily UnsolvedMath ID 2562 · Group Theory}
\subsection{Definitions and mathematical statement}

\paragraph{Shared notation.}$\mathbb N=\{1,2,\ldots\}$, and $\mathbb Z,\mathbb Q,\mathbb R,\mathbb C$ denote the integers, rationals, reals and complex numbers. Any other conventions are specified locally.


Let $G$ be a finite nonabelian simple group and $D$ a conjugacy class of involutions (elements of order $2$). Let $\Gamma$ be the complete graph with vertex set $D$, whose edge $\{a,b\}$ is coloured by the order $|ab|$ of the group element $ab$. An automorphism of this coloured graph preserves each colour. Write
\[
 \operatorname{Aut}_t(\Gamma)=\{\tau\in\operatorname{Sym}(D): |\tau(a)\tau(b)|=t\text{ whenever }a,b\in D,\ |ab|=t\}.
\]
Because $D$ is finite, this one-way preservation of a given colour is equivalent to preserving its colour class as a set, and $\operatorname{Aut}(\Gamma)=\bigcap_t\operatorname{Aut}_t(\Gamma)$. Let $p$ be the second-smallest prime divisor of $|G|$, the smallest being $2$.
\paragraph{Statement recorded in the supplied source.}
For every such group $G$ and class $D$, is
\[
 \operatorname{Aut}(\Gamma)=\operatorname{Aut}_2(\Gamma)\cap\operatorname{Aut}_p(\Gamma)?
\]
Thus, does preserving the two edge colours given by the smallest prime divisors force preservation of every product-order colour? \sref{2562}{2}
\subsection{Short English statement}
For a conjugacy class of involutions in a finite simple group, do the first two prime product orders determine every colour-preserving symmetry?
\subsection{Sources}
\begin{description}
\item[\hypertarget{source:2562:1}{S1}] ULAM.AI and the UnsolvedMath Contributors, \emph{UnsolvedMath: A Curated Collection of Open Mathematics Problems}, record 2562 (KOU-21.53). CC BY 4.0. \url{https://huggingface.co/datasets/ulamai/UnsolvedMath}.
\item[\hypertarget{source:2562:2}{S2}] I. B. Gorshkov, Problem 21.53, in E. I. Khukhro and V. D. Mazurov (eds.), \emph{The Kourovka Notebook: Unsolved Problems in Group Theory}, 21st edition, arXiv:1401.0300v48 (6 October 2026). See also source definitions in Problems 21.52. This version records a negative answer; the displayed question is the historically registered target. \url{https://arxiv.org/abs/1401.0300v48}.
\end{description}
\label{end:2562}
\end{document}
